\chapter{The Canvas and ASCII Backends}
\label{ch:canvas-ascii-backends}

These two backends share a family resemblance --- both are 2D-only, in the same spirit as
\texttt{SDL\_RENDERER} --- but arrive there by opposite routes: \texttt{CANVAS} replaces
SDL's own rendering entirely with a browser API, while \texttt{ASCII} is, quite literally, a
thin decorator wrapped around \texttt{SDL\_RENDERER}'s own real implementation.

\section{CANVAS: real code, an honest verification ceiling}

\texttt{CANVAS} is Emscripten-only, targeting the HTML \texttt{CanvasRenderingContext2D} API
(\texttt{drawImage} / \texttt{putImageData} / \texttt{fillRect}) instead of SDL3's own 2D blit
API. The single fact that qualifies every other claim in this section: this project's own
development environment has \emph{no real browser and no real DOM}. The \texttt{node}
runtime that executes this backend's compiled test suite provides no
\texttt{CanvasRenderingContext2D} at all, and \texttt{SDL\_Init(SDL\_INIT\_VIDEO)} itself
throws under Emscripten running in \texttt{node} --- confirmed directly, not assumed. Every
checkmark in this backend's own documentation therefore means ``implemented, code-reviewed,
and passing structural (non-pixel) test coverage,'' not ``pixel-verified in a real browser''
--- a narrower claim than every backend in the previous chapters, and stated as such rather
than blurred.

\subsection{Five real bugs found in an external review, before any real browser existed to test in}

An external code review found and fixed five distinct bugs purely by reasoning about Canvas
2D compositing semantics against the CSS Compositing specification, without ever running the
code in a browser:

\begin{enumerate}
  \item \cnaclass{BlendState.Opaque} cleared the \emph{entire} canvas, not just the drawn
        sprite --- a direct consequence of using \texttt{globalCompositeOperation = 'copy'}
        with no clip region, which wipes everything outside the shape actually drawn, per the
        Porter-Duff \texttt{copy} operator's own definition. Fixed by clipping to the exact
        drawn rectangle before compositing.
  \item \texttt{AlphaBlend} and \texttt{NonPremultiplied} were treated as identical, on the
        original reasoning that this backend's own textures are never genuinely
        premultiplied --- an assumption that does not hold project-wide, since
        \texttt{SDL\_RENDERER}'s own test suite deliberately constructs premultiplied test
        data. Fixed with a real per-pixel un-premultiply pass applied only to the
        \texttt{AlphaBlend} case.
  \item RGB tint visibly darkened or lightened semi-transparent edges, traced to an
        alpha-squared error in the original multiply-plus-destination-in compositing trick,
        confirmed wrong against the CSS Compositing specification's own blend formula. Fixed
        with a direct per-pixel RGB multiply that leaves the alpha channel untouched.
  \item Invalid combinations of wrap/mirror address modes were silently skipped rather than
        raising an error --- a direct contradiction of this project's own ``throw rather than
        guess'' convention (Chapter~\ref{ch:design-philosophy}). Fixed by moving the
        validation into C\texttt{++} itself, throwing a real \texttt{std::runtime\_error}
        before the call ever reaches JavaScript at all.
  \item A cached mirrored-tile pattern was never invalidated after a pixel update --- fixed
        in two separate passes, first for an explicit \texttt{Texture2D::SetData} call, then
        again for the case of a bound \cnaclass{RenderTarget2D} whose pixels change through
        direct draws while still bound, a path that never goes through \texttt{SetData} at
        all and so was missed by the first fix.
\end{enumerate}

A closely related sixth fix: \texttt{Clear()} was blending against old canvas content under a
leftover composite mode, and permanently reset the canvas transform rather than
saving/restoring it --- both corrected together.

\subsection{Not-fully-native features, honestly labeled}

\cnaclass{TextureAddressMode::Wrap} is implemented via \texttt{ctx.createPattern(source,
'repeat')}, applied only when the source rectangle actually exceeds the texture's bounds;
\texttt{Mirror} uses a lazily-built, cached $2\times2$ pre-tiled mirrored canvas. Both carry
the same real-browser verification caveat as everything else in this section. Mixed per-axis
combinations (wrap on one axis, mirror or clamp on another) throw by design, per the fix
above.

\subsection{A worked example: the un-premultiply pass, in real numbers}

The \texttt{AlphaBlend} / \texttt{NonPremultiplied} fix above (bug~2) is worth the same
hand-computed treatment the ASCII backend's own quantizer gets below, since the actual
per-pixel formula --- read directly from \texttt{CanvasSpriteBatchBackend.cpp}'s own
Canvas-2D-side pixel loop --- is small enough to verify by hand. A half-transparent, dark-red
premultiplied source pixel, $\text{RGBA} = (64, 32, 16, 128)$ (alpha $\approx 50\%$), is the
kind of value \texttt{SDL\_RENDERER}'s own test suite deliberately constructs, per the bug's
own root-cause description above:

\begin{lstlisting}
const inv = 255 / aa;              // aa = 128 here -> inv = 1.9921875
rr = Math.min(255, rr * inv);      // 64 * 1.9921875 = 127.5
gg = Math.min(255, gg * inv);      // 32 * 1.9921875 = 63.75
bb = Math.min(255, bb * inv);      // 16 * 1.9921875 = 31.875
// alpha (aa = 128) is read and written back completely untouched;
// data[] is a Uint8ClampedArray, so each float above is rounded/clamped to a byte on assignment
\end{lstlisting}

The un-premultiplied result, $\approx(128, 64, 32, 128)$, is close to double the original
channel values --- the pixel's true, alpha-independent color, exactly what a
\texttt{NonPremultiplied}-style source already stores directly, and exactly
why treating the two blend modes as identical (the original, buggy assumption) silently
darkened every partially-transparent \texttt{AlphaBlend} pixel by roughly the same factor as
its own alpha: without the \texttt{inv} multiply, the raw premultiplied RGB $(64,32,16)$ would
have been drawn as though it were the pixel's real, full-brightness color, when the true color
is nearly twice as bright in each channel. \texttt{Math.min(255, \ldots{})} is not a decorative
safety clamp either --- it is load-bearing for any source alpha below 128, where
\texttt{inv} exceeds 2 and an already-bright channel value can genuinely overflow past 255
before clamping.

\subsection{What real browser verification would still require}

The backend's own documentation includes a full thirteen-item manual browser verification
checklist --- explicitly left entirely unchecked, described as ``the closest honest
equivalent to SDL\_RENDERER's own pixel-verified completeness table,'' pending a human
running it in a real browser (for instance via Emscripten's \texttt{emrun}).

\section{ASCII: composition, not reimplementation}

\texttt{ASCII} is not a real terminal or TTY backend --- it renders inside an ordinary SDL
window, exactly like \texttt{SDL\_RENDERER}, and is architecturally, in the project's own
words, ``a thin decorator around SDL\_RENDERER's own \texttt{SdlGraphicsBackend}.'' A game
draws normally into a private offscreen target; \texttt{Present()} reads that frame back,
quantizes it into a grid of glyph-plus-color cells, and draws the resulting grid onto the
real window using the internal \cnaclass{ISpriteBatchBackend} / \cnaclass{ITextureBackend}
pair directly --- deliberately \emph{not} going through the public \cnaclass{SpriteFont} API,
consistent with the project's own rule that backend internals stay below, and hidden from,
the public XNA-facing API surface.

\subsection{A genuinely distinctive claim: no blocked decisions, no environment gate}

This backend's own completeness documentation makes a claim none of the others in this
chapter or the previous ones can: ``no BLOCKED section and no \texttt{needs\_human} gate
anywhere --- every row below is automated and pixel-verified in this dev environment\ldots
including real-window \texttt{Present()} output.'' Because it inherits \texttt{SDL\_RENDERER}'s
own implementation wholesale for everything a game actually draws --- \cnaclass{SpriteBatch},
\cnaclass{Texture2D}, \cnaclass{SpriteFont}, \cnaclass{BlendState}, \cnaclass{SamplerState},
render targets --- it also inherits that backend's entire, exhaustively-verified completeness
table transparently, without needing to repeat it. The one partial exception:
\texttt{SetRenderTarget(nullptr)} is intercepted and redirected to the private game-target
buffer rather than the literal window backbuffer, ``so the game can never accidentally draw
straight onto the real window.''

\subsection{The quantizer, and a deliberately hand-authored font}

Glyph selection works by standard luminance ranking (the usual $0.299R + 0.587G + 0.114B$
formula), mapped onto a ten-character density ramp (\texttt{" .:-=+*\#\%@"}) by linear rank ---
verified to be strictly increasing in visual density. Two runtime modes,
\texttt{BLACKWHITE} and \texttt{COLOR}, are selected via the \texttt{CNA\_ASCII\_MODE}
environment variable, with a runtime-configurable cell size (default $8\times8$) via
\texttt{SetCellSize()} / \texttt{GetCellSize()}. The glyph font itself is a deliberate,
hand-authored $8\times8$ pixel bitmap for exactly the ten ramp characters --- explicitly not
a vendored font file or classic code-page glyph set, a choice made specifically to avoid
introducing any new asset or license dependency at all. A dirty-cell diffing optimization is
explicitly not implemented, and explicitly not treated as a correctness gap either --- left
for later only if performance ever actually demands it.

\subsection{A worked example: from one pixel to one glyph}

The quantizer's own formula and ramp are precise enough to walk through by hand for a single
pixel, which makes the mapping concrete rather than only described:

\begin{lstlisting}
// A mid-brightness orange pixel, e.g. sampled from a sprite:
Color pixel(230, 126, 34, 255);

float luminance = 0.299f * pixel.R + 0.587f * pixel.G + 0.114f * pixel.B;
// = 0.299*230 + 0.587*126 + 0.114*34 = 68.77 + 73.96 + 3.88 = 146.6

const char* ramp = " .:-=+*#%@";       // 10 characters, index 0 = darkest
int rampIndex = static_cast<int>(luminance / 255.0f * 9.0f); // = 5 -> '+'
char glyph = ramp[rampIndex];
\end{lstlisting}

A luminance of 146.6 out of 255 (roughly 57\%) lands on \texttt{'+'}, the sixth character in
the ten-character ramp --- not the visual midpoint of the string by character count, but by
the ramp's own \emph{verified-strictly-increasing} visual density, which is exactly the
property the chapter's own verification claim rests on. A pure white pixel
(\texttt{luminance=255}) maps to index 9, \texttt{'@'}, the densest glyph; pure black maps to
index 0, a plain space --- background, rendered as the visual absence of any glyph at all,
which is also precisely why the \texttt{BlendState.AlphaBlend} bug above was so easy to miss
purely by inspection: a background that fails to alpha-blend and an intentionally-empty
space-character cell can look identical in a screenshot.

\subsection{One real, found-and-fixed bug}

Background fills plus glyph tinting never actually went through
\texttt{GraphicsDevice.BlendState} at all, so the underlying SDL renderer's blend mode
defaulted to a non-alpha-aware mode --- background fills were effectively invisible, because
an ``off'' glyph pixel simply overwrote the background with whatever raw black RGB value it
stored. Fixed by forcing real \cnaclass{BlendState.AlphaBlend} factors before every grid draw,
independent of whatever \cnaclass{BlendState} the game itself had set.

\subsection{A methodology finding, not a bug}

Reading pixels back immediately after a real \texttt{Present()} call returns stale or
garbage content, because \texttt{SDL\_RENDERER}/OpenGL present via a genuine double-buffer
swap, not a copy --- there is nothing left to read where the just-presented frame was. The
fix is test-only: a \texttt{DrawQuantizedGridForTesting()} entry point that skips the swap,
paired with a \texttt{Read\allowbreak{}RealBackbuffer\allowbreak{}ForTesting()} accessor --- real game code always goes
through the genuine \texttt{Present()} path, unaffected by either.

\subsection{Input, unmodified}

Input on this backend required zero new code at all --- verified twice, independently: once
through the existing backend-agnostic mouse test suite, and once through a new, dedicated
input test written specifically for this backend, passing in full.
